\documentclass[12pt]{article}
\usepackage[utf8]{inputenc}
\usepackage[russian]{babel}
\usepackage{amsmath}
\usepackage{amssymb}
\usepackage{geometry}
\geometry{a4paper, margin=2cm}

\begin{document}

\section*{Задача 3.1.2}

\textbf{Условие:} Вывести числа в последовательности \(a_1, \dots, a_n\), которые делятся без остатка на \(d\) (0.5 балла):\\
in: d, n, a1, ..., an \\
out: b1, ..., bk

\section*{Решение}

\begin{verbatim}
Read 2  ; d
Read 3  ; n

WHILE: Load 3
    JZero END
    Read 0
    Store 1 ; r1 := a_n
    Div 2
    Mult 2
    Sub 1
    JZero WRITE
    Jump DEC_WHILE
WRITE: Write 1
DEC_WHILE: Load 3
    Sub =1
    Store 3 ; r3--
    Jump WHILE
END: Halt 0
\end{verbatim}

\section*{Сложность}

Худшего случая как такового нет, так как всегда программа работает за одно и то же число шагов. Например, если число делится на d, то происходит переход на WRITE с одной командой, а затем выполняется DEC\_WHILE. Если число не делится, то происходит шаг на команду Jump DEC\_WHILE, а затем выполняется DEC\_WHILE. Ну и так как критерий равномерный, то их сложность совпадает. Далее для определенности рассмотрим случай, когда все числа делятся на d.

\begin{enumerate}
    \item Команды до WHILE выполняются за 2 шага
    \item Цикл выполняется n раз
    \item Тело цикла выполняется за 8 шагов
    \item Write 1 за 1 шаг
    \item Декрементация за 4 шага
    \item Последняя проверка условия (когда n==0) и остановка занимают 3 шага
\end{enumerate}

Значит, сложность в худшем случае:
\[
    T_{\max} = 13n + 5
\]

\end{document}